\cleardoublepage
\chapter{Penutup}
\label{chap:penutup}

\section{Rangkuman Pelaksanaan Pekerjaan}
Pada bulan Juli dan Agustus 2026, penulis melaksanakan kerja praktik di PT Nutech Integrasi sebagai \textit{Product Specialist Intern} pada Divisi \textit{Business \& Product Development}. Kegiatan bulan Juli meliputi pengenalan portofolio perusahaan, pemeriksaan katalog produk, pemahaman arsitektur \textit{Automatic Fare Collection} (AFC), serta eksplorasi \textit{Smart Depot Management System} (SDMS) dan pengembangan prototipe \textit{Operational Dashboard}. Sementara itu, kegiatan bulan Agustus mencakup pembahasan solusi Penerangan Jalan Umum Tenaga Surya (PJUTS) bersama PT Jarwinn Solar Panel, kunjungan ke pameran \textit{Digital Transformation Indonesia} (DTI), diskusi internal divisi, serta presentasi akhir kerja praktik.

\section{Rangkuman Hasil Pekerjaan}
Seluruh penugasan dari mentor telah dikerjakan dan menghasilkan hasil yang sesuai dengan ruang lingkup pekerjaan. Hasil yang didapat meliputi perbaikan draf katalog produk, identifikasi kebutuhan informasi monitoring depo Kereta Api Indonesia, serta prototipe \textit{Operational Dashboard} berbasis web. Prototipe ini telah menampilkan informasi, tetapi belum terintegrasi dengan sistem operasional depo. Melalui pelaksanaan kerja praktik, penulis memperoleh pemahaman mengenai kebutuhan informasi, pengembangan antarmuka web, serta peningkatan kemampuan komunikasi dan kerja sama dalam lingkungan profesional.

\section{Evaluasi Diri}
Selama berlangsungnya kerja praktik, penulis memperoleh sejumlah pengalaman dalam menerapkan tanggung jawab terhadap pekerjaan yang ditugaskan mentor, berkomunikasi secara jelas, serta menjaga dokumen rahasia milik PT Nutech Integrasi. Penulis belajar mengenai prinsip-prinsip kedisiplinan profesional, akuntabilitas dalam setiap pengerjaan tugas, efektivitas komunikasi interpersonal, serta penerapan etika dalam menjaga kerahasiaan. Salah satu tantangan yang dihadapi adalah proses adaptasi terhadap pengetahuan teknis yang spesifik di bidang perkeretaapian dan sistem Penerangan Jalan Umum Tenaga Surya (PJUTS). Akan tetapi, untuk mengatasi hambatan tersebut, penulis melakukan beberapa upaya, salah satunya dengan mempelajari materi yang relevan, meninjau kembali arahan pembimbing yang telah diberikan, serta berdiskusi bersama mentor dan rekan kerja.

\section{Saran}
Merujuk pada pengalaman penulis selama melaksanakan kerja praktik di PT Nutech Integrasi, terdapat beberapa saran yang dapat menjadi pertimbangan bagi perusahaan, yaitu:

\begin{enumerate}
    \item Pendampingan mentor sebaiknya tetap dilaksanakan agar mahasiswa kerja praktik memperoleh arahan yang jelas dan dapat memahami penugasan dengan baik.
    \item Dokumentasi produk dan teknologi diharapkan terus dikembangkan untuk memudahkan pencarian informasi teknis serta mendukung pelaksanaan pekerjaan.
    \item Kegiatan diskusi dan presentasi internal yang dilaksanakan secara berkala sebaiknya dipertahankan sebagai sarana untuk menyampaikan perkembangan pekerjaan, melatih keterampilan presentasi, meningkatkan keberanian dalam mengutarakan pendapat, serta mengembangkan kemampuan komunikasi.
\end{enumerate}

\section{Ucapan Terima Kasih}
Penulis menyampaikan terima kasih kepada keluarga, Program Studi Sistem dan Teknologi Informasi Institut Teknologi Bandung, dosen pembimbing akademik, serta koordinator kerja praktik atas dukungan dan arahan yang diberikan. Penulis juga berterima kasih kepada PT Nutech Integrasi, khususnya Divisi \textit{Business \& Product Development}, Kak Hani selaku mentor, Mas Reza, serta seluruh rekan kerja dan karyawan yang memberikan bimbingan, bantuan, dan pengalaman berharga selama pelaksanaan kerja praktik.
